% BHU PYQ -- Manually transcribed & error-corrected
% Paper: BCF-116 : Business Statistics
% Exam: B.Com. (Hons.) FMM Semester I Examination, 2023-24

\documentclass[11pt,a4paper]{article}
\usepackage[a4paper,top=2.5cm,bottom=2.5cm,left=2.8cm,right=2.8cm,headheight=14pt]{geometry}
\usepackage{amsmath,amssymb}
\usepackage[T1]{fontenc}
\usepackage[utf8]{inputenc}
\usepackage{lmodern,microtype,fancyhdr,enumitem,hyperref,lastpage,parskip,booktabs}

\hypersetup{
    colorlinks=true,
    urlcolor=black,
    linkcolor=black,
    pdftitle={BCF-116: Business Statistics -- BHU B.Com. (Hons.) FMM Sem I 2023-24}
}

\pagestyle{fancy}
\fancyhf{}
\renewcommand{\headrulewidth}{0.4pt}
\renewcommand{\footrulewidth}{0.4pt}
\fancyhead[L]{\small   \textit{Banaras Hindu University}}
\fancyhead[R]{\small   \textit{BCF-116: Business Statistics}}
\fancyfoot[C]{\small Page  \thepage\ of \pageref{LastPage}}


\newlist{parts}{enumerate}{2}
\setlist[parts,1]{label=(\alph*),leftmargin=*,itemsep=5pt,topsep=3pt}
\setlist[parts,2]{label=(\roman*),leftmargin=*,itemsep=3pt,topsep=2pt}

\newcommand{\pts}[1]{\hfill   \textit{[#1]}}


\begin{document}


\begin{center}
  {\large\bfseries BANARAS HINDU UNIVERSITY}\\[2pt]
  {\normalsize B.Com. (Hons.) FMM --- Semester I Examination, 2023-24}\\[6pt]
  \rule{0.8\linewidth}{0.5pt}\\[6pt]
  {\large\bfseries Paper No. BCF-116: Business Statistics}\\[4pt]
  {\normalsize Time: 3 Hours\hfill Maximum Marks: 70}
\end{center}

\rule{\linewidth}{0.4pt}

\smallskip



\noindent   \textit{Instructions: Answer all questions. All questions carry equal marks (14 marks each).}

\bigskip




\noindent   \textbf{Question 1.}
What do you mean by Primary and Secondary Data? What methods are used to collect these? Explain their relative merits and demerits.\pts{14}

\centerline{   \textbf{OR}}

Rearrange the following series with equal intervals and then prepare 'less than' and 'more than' cumulative frequency distribution:\pts{14}

\begin{center}
  \small
  
\begin{tabular}{cc}
     oprule
   \textbf{Class Interval} &   \textbf{Frequency} \\
    \midrule
    0--5   & 3   \\
    5--6   & 2   \\
    6--9   & 7   \\
    9--12  & 5   \\
    12--17 & 16  \\
    17--18 & 12  \\
    18--20 & 15  \\
    20--24 & 20  \\
    24--25 & 8   \\
    25--30 & 10  \\
    30--36 & 2   \\
    \midrule
   \textbf{Total} &   \textbf{100} \\
    ottomrule
  \end{tabular}
\end{center}

\medskip\rule{\linewidth}{0.2pt}




\noindent   \textbf{Question 2.}
Explain what is meant by 'Central Tendency' and describe the various methods of measuring it.\pts{14}

\centerline{   \textbf{OR}}

Calculate the mean and median from the following series:\pts{14}

\begin{center}
  \small
  
\begin{tabular}{cc}
     oprule
   \textbf{Class Interval} &   \textbf{Frequency} \\
    \midrule
    0--9   & 32 \\
    10--19 & 36 \\
    20--29 & 20 \\
    30--39 & 30 \\
    40--59 & 48 \\
    60--79 & 24 \\
    80--99 & 2  \\
    ottomrule
  \end{tabular}
\end{center}

\medskip\rule{\linewidth}{0.2pt}

\pagebreak




\noindent   \textbf{Question 3.}
What do you understand by Skewness? Point out its role in analyzing a frequency distribution.\pts{14}

\centerline{   \textbf{OR}}

Calculate the standard deviation and its coefficient from the following data:\pts{14}

\begin{center}
  \small
  
\begin{tabular}{cc}
     oprule
   \textbf{Marks} &   \textbf{Frequency} \\
    \midrule
    10--20 & 6  \\
    20--30 & 8  \\
    30--40 & 15 \\
    40--50 & 12 \\
    50--60 & 6  \\
    60--70 & 2  \\
    70--80 & 1  \\
    ottomrule
  \end{tabular}
\end{center}

\medskip\rule{\linewidth}{0.2pt}




\noindent   \textbf{Question 4.}
What are the components of a time series? How would you find out the trend values in a time series by the method of least squares?\pts{14}

\centerline{   \textbf{OR}}

Find out the value of the trend by three-yearly moving averages and plot them on a graph paper using the following data:\pts{14}

\begin{center}
  \small
  
\begin{tabular}{lccccccccccccc}
     oprule
   \textbf{Year}  & 1995 & 1996 & 1997 & 1998 & 1999 & 2000 & 2001 & 2002 & 2003 & 2004 & 2005 & 2006 & 2007 \\
    \midrule
   \textbf{Value} & 10   & 15   & 12   & 18   & 15   & 22   & 19   & 24   & 20   & 26   & 22   & 30   & 25   \\
    ottomrule
  \end{tabular}
\end{center}

\medskip\rule{\linewidth}{0.2pt}




\noindent   \textbf{Question 5.}
How are index numbers constructed? Discuss the importance of the choice of base year and selection of weights in the construction of cost of living Index Numbers.\pts{14}

\centerline{   \textbf{OR}}

Calculate Fisher's Ideal Index number from the data given below:\pts{14}

\begin{center}
  \small
  
\begin{tabular}{ccccc}
     oprule
   \textbf{Item} &   \textbf{Base Price} ($p_0$) &   \textbf{Base Quantity} ($q_0$) &   \textbf{Current Price} ($p_1$) &   \textbf{Current Quantity} ($q_1$) \\
    \midrule
    A & 6  & 50  & 10 & 60  \\
    B & 2  & 100 & 2  & 120 \\
    C & 4  & 60  & 6  & 60  \\
    D & 12 & 40  & 15 & 25  \\
    E & 8  & 50  & 12 & 35  \\
    ottomrule
  \end{tabular}
\end{center}

\vfill
\rule{\linewidth}{0.4pt}

\begin{center}\small
\href{https://bhu-exam.vercel.app/}{Click here to visit our website}\\[4pt]
\href{https://me-aryan.vercel.app/}{Click here to visit our contributor}\\[4pt]
\href{https://quashan-me.vercel.app/}{Click here to visit our contributor}
\end{center}

\end{document}
